\chapter{Release 2 : Collecte intelligente et IoT}
\label{chap:release2}

% Affichage robuste des diagrammes : le chapitre reste compilable
% même lorsque les PNG n'ont pas encore été exportés depuis Draw.io.
\providecommand{\releasefigure}[4]{%
\releasefigure{image/usecase_release2_collecte_intelligente.png}{0.98}{Diagramme de cas d'utilisation -- Release 2}{usecase_release2_collecte_intelligente.drawio}
\label{fig:r2_usecase}


\subsection{Diagrammes de séquence du Sprint 3}

Deux scénarios principaux sont retenus pour représenter les interactions du Sprint 3.

\subsubsection{Scénario 1 : dépôt des déchets et scan du QR code}

Le citoyen présente son QR code au dispositif. Le système le scanne, vérifie son identité et son rôle, puis autorise ou refuse la poursuite de l'opération. Lorsque le QR code est valide, le citoyen sélectionne le type de déchet, effectue le dépôt et le système traite les informations correspondantes.

Le diagramme utilise :
\begin{itemize}
    \item \texttt{alt} pour distinguer un QR code valide d'un QR code invalide ;
    \item \texttt{ref} pour représenter la sous-interaction d'attribution des points ;
    \item \texttt{loop} pour représenter la répétition des dépôts.
\end{itemize}

\releasefigure{image/sequence_release2_sprint3_depot_citoyen.png}{0.98}{Diagramme de séquence -- Dépôt des déchets et scan du QR code}{sequence_release2_sprint3_depot_citoyen.drawio}
\label{fig:r2_s3_seq_depot}


\subsubsection{Scénario 2 : types de déchets et points de collecte}

Le citoyen consulte les points de collecte disponibles. Le système récupère les informations enregistrées et les affiche. Pour chaque point, le type de déchets accepté peut être vérifié afin d'orienter le citoyen vers un point adapté.

Le fragment \texttt{loop} représente la consultation successive des points de collecte. Le fragment \texttt{alt} distingue les types de déchets acceptés et non acceptés. Le fragment \texttt{ref} représente la mise à jour des informations relatives au point de collecte.

\releasefigure{image/sequence_release2_sprint3_vidage_collecteur.png}{0.98}{Diagramme de séquence -- Types de déchets et points de collecte}{sequence_release2_sprint3_vidage_collecteur.drawio}
\label{fig:r2_s3_seq_points}


\section{Sprint 4 : Gestion opérationnelle et SmartBin}
\label{sec:r2_sprint4}

\subsection{Objectifs du Sprint 4}

Le Sprint 4, d'une durée de \textbf{12 jours}, complète le Sprint 3 par les fonctionnalités de gestion opérationnelle et par l'intégration de la poubelle intelligente.

Il couvre cinq éléments principaux : la gestion des missions de collecte, la mise à jour de leur état, le suivi des centres et des équipements, l'intégration du SmartBin et la transmission des données de remplissage. L'administrateur dispose également d'une vue permettant de suivre l'état des poubelles.

\subsection{Backlog du Sprint 4}

\begin{longtable}{|>{\centering\arraybackslash}p{1cm}|p{9.3cm}|>{\centering\arraybackslash}p{2cm}|}
\caption{Backlog du Sprint 4}
\label{tab:r2_s4_backlog}
\\ \hline
\rowcolor{gray!25}
\textbf{ID} & \textbf{User Story} & \textbf{Priorité}
\\ \hline
\endfirsthead
\hline
\rowcolor{gray!25}
\textbf{ID} & \textbf{User Story} & \textbf{Priorité}
\\ \hline
S4.1 & En tant que collecteur, je veux consulter mes missions de collecte et mettre à jour leur état. & Élevée \\ \hline
S4.2 & En tant qu'administrateur, je veux suivre les centres de collecte et leurs équipements. & Élevée \\ \hline
S4.3 & En tant que système, je veux intégrer les SmartBins à la plateforme EcoRewind. & Élevée \\ \hline
S4.4 & En tant que système, je veux transmettre les données de remplissage des SmartBins. & Élevée \\ \hline
S4.5 & En tant qu'administrateur, je veux suivre l'état et le niveau de remplissage des poubelles. & Élevée \\ \hline
\end{longtable}

\subsection{Missions de collecte}

Une mission de collecte représente une opération attribuée à un collecteur. Elle possède un état permettant de suivre son avancement, par exemple \textit{en attente}, \textit{en cours} ou \textit{terminée}.

Lorsqu'une mission est sélectionnée, le collecteur peut faire évoluer son état. Le système enregistre la modification afin que les autres acteurs disposent d'une information actualisée.

\subsection{Suivi des centres et des équipements}

Les centres de collecte constituent les points organisationnels auxquels peuvent être associés plusieurs équipements. L'administrateur peut consulter les centres ainsi que l'état de leurs équipements.

Cette fonctionnalité permet de centraliser les informations nécessaires au suivi opérationnel et d'identifier les équipements nécessitant une intervention.

\subsection{Intégration de la poubelle intelligente}

La SmartBin constitue le composant IoT de la Release 2. Elle collecte des informations provenant des capteurs puis les transmet à la plateforme. Les informations principales associées à une poubelle comprennent son identifiant, sa localisation, son type de déchet, son poids et son état.

La communication entre le prototype et la plateforme est réalisée par l'intermédiaire du réseau Wi-Fi et d'un service de données distant. Les informations reçues peuvent ensuite être exploitées par l'application pour assurer le suivi des poubelles.

\subsection{Transmission des données de remplissage}

Le niveau de remplissage est déterminé à partir des données de poids du SmartBin. Les mesures sont transmises à la plateforme afin de conserver un état actualisé de chaque poubelle.

Lorsque le niveau atteint une condition nécessitant une intervention, le système peut identifier la poubelle comme prioritaire et transmettre une notification aux acteurs concernés.

\subsection{Suivi des poubelles par l'administrateur}

L'administrateur dispose d'une vue de suivi regroupant les informations des SmartBins. Il peut notamment consulter leur état, leur niveau de remplissage et leur localisation.

Cette vue permet de rapprocher les données provenant du dispositif IoT des informations opérationnelles relatives aux centres et aux missions de collecte.

\subsection{Diagramme de classes}

Le diagramme de classes de la Release 2 regroupe les principales entités nécessaires aux deux sprints. Le modèle comprend notamment \texttt{User}, \texttt{BinScan}, \texttt{CollectionPoint}, \texttt{MissionCollecte}, \texttt{CentreCollecte}, \texttt{Equipement}, \texttt{SmartBin}, \texttt{DonneeRemplissage} et \texttt{CollectorLog}.

Les relations entre ces classes permettent de relier les utilisateurs aux opérations de dépôt, les centres aux équipements et aux poubelles, ainsi que les missions aux opérations de collecte.

\releasefigure{image/class_release2_collecte_intelligente.png}{0.98}{Diagramme de classes -- Release 2}{class_release2_collecte_intelligente.drawio}
\label{fig:r2_class}


\subsection{Diagrammes de séquence du Sprint 4}

\subsubsection{Scénario 1 : mission de collecte}

Le collecteur consulte les missions disponibles. Le système retourne les missions ainsi que leur état. Lorsqu'une mission est sélectionnée, son état est mis à jour puis enregistré.

Le fragment \texttt{alt} distingue notamment une mission en attente d'une mission déjà terminée. Le fragment \texttt{loop} permet de représenter le traitement successif des missions d'une tournée. Le fragment \texttt{ref} représente l'enregistrement de l'opération de collecte.

\releasefigure{image/sequence_release2_sprint4_missions_collecte.png}{0.98}{Diagramme de séquence -- Missions de collecte et mise à jour de leur état}{sequence_release2_sprint4_missions_collecte.drawio}
\label{fig:r2_s4_seq_missions}


\subsubsection{Scénario 2 : transmission du remplissage et suivi du SmartBin}

Le SmartBin mesure les données nécessaires au calcul du remplissage et les transmet à EcoRewind. Les données sont enregistrées afin de mettre à jour l'état de la poubelle.

Le fragment \texttt{alt} distingue le cas où le niveau de remplissage nécessite une intervention du cas où le niveau reste normal. Le fragment \texttt{loop} représente la surveillance périodique des poubelles. Le fragment \texttt{ref} représente la consultation des SmartBins par l'administrateur.

\releasefigure{image/sequence_release2_sprint4_smartbin_remplissage.png}{0.98}{Diagramme de séquence -- Transmission des données de remplissage et suivi du SmartBin}{sequence_release2_sprint4_smartbin_remplissage.drawio}
\label{fig:r2_s4_seq_smartbin}


\section{Bilan de la Release 2}

La Release 2 est ainsi structurée de manière progressive. Le Sprint 3, réalisé sur 11 jours, met en place le parcours fonctionnel du dépôt : identification par QR code, gestion des types de déchets et consultation des points de collecte. Le Sprint 4, réalisé sur 12 jours, ajoute la dimension opérationnelle : missions de collecte, suivi des centres et équipements, intégration des SmartBins, transmission des données de remplissage et suivi des poubelles par l'administrateur.

Cette organisation permet de séparer clairement les fonctionnalités orientées vers le citoyen des fonctionnalités de supervision et de gestion de la collecte, tout en préparant leur interaction au sein de la même plateforme EcoRewind.

\section{Conclusion}

La Release 2 apporte à EcoRewind une chaîne fonctionnelle complète autour de la collecte intelligente. Le Sprint 3 établit le parcours de dépôt et les informations nécessaires à l'orientation du citoyen. Le Sprint 4 étend ce fonctionnement aux acteurs opérationnels et au dispositif IoT.

Les diagrammes de cas d'utilisation, de classes et de séquence permettent de représenter cette évolution à différents niveaux : les fonctionnalités attendues, la structure des données et les interactions entre les acteurs, la plateforme et les SmartBins.
